%% paper47_representation_theory_en.tex
%% Paper 47: Representation Theory — From Group Algebras to the Langlands Program
%% Author: Michael Fuhrmann
%% Project: Specialist Mathematics, Build 135
%% Date: 2026-03-12

\documentclass[12pt,a4paper]{amsart}
\usepackage{amsmath,amssymb,amsthm}
\usepackage{mathtools}
\usepackage[utf8]{inputenc}
\usepackage[T1]{fontenc}
\usepackage{hyperref}
\usepackage{enumitem,booktabs,array}
\usepackage{geometry}
\geometry{margin=2.5cm}

\newtheorem{theorem}{Theorem}[section]
\newtheorem{lemma}[theorem]{Lemma}
\newtheorem{corollary}[theorem]{Corollary}
\newtheorem{proposition}[theorem]{Proposition}
\newtheorem{conjecture}[theorem]{Conjecture}
\theoremstyle{definition}
\newtheorem{definition}[theorem]{Definition}
\newtheorem{example}[theorem]{Example}
\newtheorem{remark}[theorem]{Remark}

\title{Representation Theory:\\
From Group Algebras to the Langlands Program}

\author{Michael Fuhrmann}
\address{Specialist Mathematics Project, Build 135}
\email{specialist-maths@localhost}
\date{\today}

\begin{document}

\maketitle

\begin{abstract}
Representation theory provides a systematic framework for understanding abstract
groups through their concrete actions on vector spaces.  Beginning with the
fundamental definitions of group homomorphisms into general linear groups, we
develop the core machinery of the subject: Schur's Lemma, Maschke's Theorem on
complete reducibility, and the rich orthogonality theory of characters.  We then
examine the representation theory of the symmetric groups $S_n$ via Young tableaux
and Specht modules, and pass to the continuous setting with the highest-weight
classification of representations of $\mathrm{GL}_n$ and $\mathrm{SL}_n$.
Applications range from Burnside's $pq$-theorem in pure group theory to molecular
symmetry and crystallography in chemistry and physics.  We conclude with an
introduction to Galois representations and the Langlands correspondence, one of
the deepest conjectural frameworks in modern mathematics.
\end{abstract}

\tableofcontents

%% ============================================================
\section{Introduction}
%% ============================================================

Representation theory is the art of understanding abstract algebraic structures
by letting them act on vector spaces.  A \emph{representation} of a group $G$ on
a vector space $V$ is simply a group homomorphism
\[
    \rho \colon G \longrightarrow \mathrm{GL}(V),
\]
where $\mathrm{GL}(V)$ denotes the group of invertible linear endomorphisms of
$V$.  By choosing a basis one may identify $\mathrm{GL}(V) \cong \mathrm{GL}(n,k)$
for $n = \dim_k V$, so that $\rho$ associates to each group element $g \in G$ an
invertible $n \times n$ matrix $\rho(g)$.

The subject sits at a remarkable crossroads of algebra and linear algebra.
On the one hand it interprets the often intractable combinatorics of a group in
terms of the well-understood geometry of matrix multiplication.  On the other hand
it provides the algebraic backbone for much of modern physics: the representations
of the rotation group $\mathrm{SO}(3)$ and its double cover $\mathrm{SU}(2)$
describe quantum-mechanical spin, while the representations of the Lorentz group
and $\mathrm{SU}(3)$ classify elementary particles in the Standard Model.

From a purely mathematical perspective the subject has driven enormous advances.
Burnside's 1904 proof that every group of order $p^a q^b$ ($p$, $q$ prime) is
solvable relies on character theory in an essential way that has resisted any
purely group-theoretic proof to this day.  The classification of finite simple
groups could not have been accomplished without the detailed study of modular
representations.  And at the frontier, the Langlands program---formulated by
Robert Langlands in a famous 1967 letter to André Weil---predicts deep
correspondences between automorphic representations of reductive groups and Galois
representations, unifying vast swaths of number theory, geometry, and analysis.

\medskip

\noindent\textbf{Overview.}  Section~\ref{sec:basic} introduces the basic
definitions and the most important classes of examples.  Section~\ref{sec:schur}
proves Schur's Lemma and draws its fundamental consequences.
Section~\ref{sec:maschke} establishes complete reducibility over fields whose
characteristic does not divide the group order.  Section~\ref{sec:characters}
develops character theory including the two orthogonality relations.
Section~\ref{sec:sn} treats the representation theory of the symmetric group via
Young tableaux.  Section~\ref{sec:gln} covers the highest-weight theory for
$\mathrm{GL}_n$ and $\mathrm{SL}_n$.  Section~\ref{sec:applications} collects
three classical applications.  Section~\ref{sec:langlands} gives an introduction
to Galois representations and the Langlands program.

%% ============================================================
\section{Group Representations}
\label{sec:basic}
%% ============================================================

Throughout let $G$ be a group and $k$ a field.  Unless otherwise stated we take
$k = \mathbb{C}$.

\begin{definition}
A \emph{representation} of $G$ over $k$ is a pair $(V,\rho)$ where $V$ is a
$k$-vector space and $\rho \colon G \to \mathrm{GL}(V)$ is a group homomorphism.
The \emph{dimension} (or \emph{degree}) of $(V,\rho)$ is $\dim_k V$.
\end{definition}

Equivalently a representation is a $k[G]$-module, where $k[G]$ denotes the
\emph{group algebra} of $G$ over $k$: the free $k$-vector space on the elements
of $G$ with multiplication extending that of $G$ by linearity.

\begin{example}[Trivial representation]
The \emph{trivial representation} is $\rho_{\mathrm{triv}} \colon G \to
\mathrm{GL}(k)$, $g \mapsto 1$.  Every element acts as the identity.
\end{example}

\begin{example}[Regular representation]
The \emph{regular representation} $(\mathbb{C}[G], \rho_{\mathrm{reg}})$ is
defined by left multiplication:
\[
    \rho_{\mathrm{reg}}(g)(e_h) = e_{gh},
\]
where $\{e_h\}_{h \in G}$ is the standard basis of $\mathbb{C}[G]$.  Its
dimension is $|G|$ and it contains every irreducible representation of $G$ with
multiplicity equal to its own dimension.
\end{example}

\begin{example}[Permutation representation]
If $G$ acts on a finite set $X = \{x_1,\ldots,x_n\}$, the \emph{permutation
representation} $V = \mathbb{C}^X$ is defined by
$\rho(g)(e_{x_i}) = e_{g \cdot x_i}$.  The regular representation is the special
case $X = G$ with the left-multiplication action.
\end{example}

\begin{example}[Sign representation of $S_n$]
The symmetric group $S_n$ has a one-dimensional representation
$\mathrm{sgn} \colon S_n \to \{+1,-1\}$ sending each permutation to its sign.
For $n \geq 2$ this is distinct from the trivial representation, providing our
first example of a non-trivial irreducible representation.
\end{example}

\begin{definition}
A \emph{subrepresentation} of $(V,\rho)$ is a $G$-stable subspace $W \subseteq V$,
i.e.\ $\rho(g)W \subseteq W$ for all $g \in G$.  The representation $(V,\rho)$ is
\emph{irreducible} if its only subrepresentations are $\{0\}$ and $V$ itself.
\end{definition}

\begin{definition}
A \emph{morphism} (or \emph{$G$-equivariant map}, or \emph{intertwiner}) between
representations $(V,\rho)$ and $(W,\sigma)$ is a linear map $T \colon V \to W$
satisfying
\[
    T \circ \rho(g) = \sigma(g) \circ T \quad \text{for all } g \in G.
\]
Two representations are \emph{isomorphic} (or \emph{equivalent}) if there exists
a bijective intertwiner between them.
\end{definition}

\begin{example}[Representations of $\mathbb{Z}/2\mathbb{Z}$]
The group $\mathbb{Z}/2\mathbb{Z} = \{0,1\}$ has exactly two irreducible complex
representations, both one-dimensional:
\begin{itemize}
  \item The trivial representation: $\rho(0) = (1)$, $\rho(1) = (1)$.
  \item The sign representation: $\rho(0) = (1)$, $\rho(1) = (-1)$.
\end{itemize}
These exhaust all irreducible representations since the number of irreducibles
equals the number of conjugacy classes, which is $2$ for an abelian group of
order $2$.
\end{example}

%% ============================================================
\section{Schur's Lemma}
\label{sec:schur}
%% ============================================================

Schur's Lemma is the cornerstone of all subsequent structure theory.  It asserts
that irreducible representations are as far from having non-trivial intertwiners
as possible.

\begin{theorem}[Schur's Lemma]
\label{thm:schur}
Let $(V,\rho)$ and $(W,\sigma)$ be irreducible representations of $G$ over a
field $k$, and let $T \colon V \to W$ be a $G$-equivariant linear map.
\begin{enumerate}[label=\textup{(\roman*)}]
  \item Either $T = 0$ or $T$ is an isomorphism.
  \item If $k$ is algebraically closed and $V = W$ (with $\rho = \sigma$), then
        $T = \lambda \cdot \mathrm{id}_V$ for some scalar $\lambda \in k$.
\end{enumerate}
\end{theorem}

\begin{proof}
(i) The kernel $\ker T$ is a subrepresentation of $V$ and the image $\mathrm{im}\,T$
is a subrepresentation of $W$.  Since $V$ is irreducible, $\ker T \in \{\{0\},V\}$.
Since $W$ is irreducible, $\mathrm{im}\, T \in \{\{0\},W\}$.  If $\ker T = V$
then $T = 0$.  Otherwise $\ker T = \{0\}$ and $\mathrm{im}\, T = W$, so $T$ is
bijective.

(ii) Since $k$ is algebraically closed, $T$ has an eigenvalue $\lambda \in k$.
The map $T - \lambda\,\mathrm{id}_V$ is again $G$-equivariant and has a non-trivial
kernel, hence by part (i) it must be the zero map, giving $T = \lambda\,\mathrm{id}_V$.
\end{proof}

\begin{corollary}
\label{cor:schur_decomp}
Let $G$ be a finite group and $k = \mathbb{C}$.  If $(V,\rho)$ and $(W,\sigma)$
are non-isomorphic irreducible representations, then every $G$-equivariant map
$T \colon V \to W$ is zero.  Consequently
\[
    \mathrm{Hom}_G(V,W) = \begin{cases} \mathbb{C} & \text{if } V \cong W, \\ 0 & \text{otherwise.} \end{cases}
\]
\end{corollary}

\begin{corollary}
Every irreducible complex representation of an abelian group is one-dimensional.
\end{corollary}

\begin{proof}
For any $h \in G$, the map $\rho(h) \colon V \to V$ is $G$-equivariant (since $G$
is abelian).  By Schur's Lemma, $\rho(h) = \lambda_h \cdot \mathrm{id}_V$.  But
then every subspace of $V$ is $G$-stable.  Irreducibility forces $\dim V = 1$.
\end{proof}

\begin{remark}
The algebraic closedness hypothesis in part (ii) is necessary.  Over $\mathbb{R}$
the two-dimensional representation of $\mathrm{SO}(2)$ by rotation matrices is
irreducible, yet its endomorphisms form a ring isomorphic to $\mathbb{C}$, not
$\mathbb{R}$.
\end{remark}

%% ============================================================
\section{Complete Reducibility: Maschke's Theorem}
\label{sec:maschke}
%% ============================================================

\begin{theorem}[Maschke's Theorem]
\label{thm:maschke}
Let $G$ be a finite group and $k$ a field with $\mathrm{char}(k) = 0$ or
$\mathrm{char}(k) \nmid |G|$.  Then every representation of $G$ over $k$ is
\emph{completely reducible}: it decomposes as a direct sum of irreducible
subrepresentations.
\end{theorem}

\begin{proof}
It suffices to show that every subrepresentation has a $G$-stable complement.
Let $(V,\rho)$ be a representation and $W \subseteq V$ a subrepresentation.
Since $V$ is a finite-dimensional vector space over $k$, there exists a linear
projection $\pi_0 \colon V \to W$ (with $\pi_0|_W = \mathrm{id}_W$), but $\pi_0$
need not be $G$-equivariant.  We \emph{average} over the group to produce a
$G$-equivariant projection:
\[
    \pi \coloneqq \frac{1}{|G|} \sum_{g \in G} \rho(g) \circ \pi_0 \circ \rho(g)^{-1}.
\]
The factor $1/|G|$ is well-defined since $|G|$ is invertible in $k$ by hypothesis.
One checks:
\begin{itemize}
  \item $\pi$ is $G$-equivariant: $\rho(h) \circ \pi = \pi \circ \rho(h)$ for all $h \in G$.
  \item $\pi|_W = \mathrm{id}_W$: for $w \in W$ and any $g$, we have $\rho(g)^{-1}w \in W$
        (since $W$ is $G$-stable), hence $\pi_0(\rho(g)^{-1}w) = \rho(g)^{-1}w$,
        and $\rho(g)\rho(g)^{-1}w = w$.
  \item In particular $\pi \circ \pi = \pi$ (idempotent).
\end{itemize}
Therefore $V = W \oplus \ker\pi$ is a decomposition into $G$-stable subspaces.
By Noetherian induction on $\dim V$, the process terminates and yields a
decomposition into irreducibles.
\end{proof}

\begin{remark}
The hypothesis $\mathrm{char}(k) \nmid |G|$ is sharp.  Over a field of
characteristic $p$ dividing $|G|$, indecomposable but reducible representations
can and do occur.  The modular representation theory of finite groups (Brauer
theory) is the study of this more subtle situation.
\end{remark}

Under the hypotheses of Maschke's Theorem every representation $(V,\rho)$
decomposes as
\[
    V \cong \bigoplus_{i} V_i^{\oplus m_i},
\]
where the $V_i$ are pairwise non-isomorphic irreducible representations and $m_i
\geq 0$ are uniquely determined non-negative integers called the
\emph{multiplicities}.  The multiplicities can be computed from the character
theory developed in the next section.

%% ============================================================
\section{Characters}
\label{sec:characters}
%% ============================================================

\begin{definition}
The \emph{character} of a representation $(V,\rho)$ is the function
\[
    \chi_V \colon G \longrightarrow k, \quad \chi_V(g) = \mathrm{tr}(\rho(g)).
\]
\end{definition}

Characters are \emph{class functions}: they are constant on conjugacy classes,
since $\mathrm{tr}(ABA^{-1}) = \mathrm{tr}(B)$.  The vector space of class
functions on $G$ is denoted $\mathrm{Cl}(G)$; it has dimension equal to the
number $r$ of conjugacy classes.

\begin{proposition}
\label{prop:char_props}
Let $(V,\rho)$ and $(W,\sigma)$ be representations.  Then:
\begin{enumerate}[label=\textup{(\roman*)}]
  \item $\chi_V(e) = \dim V$.
  \item $\chi_V(g^{-1}) = \overline{\chi_V(g)}$ for $|G| < \infty$, $k = \mathbb{C}$.
  \item $\chi_{V \oplus W} = \chi_V + \chi_W$.
  \item $\chi_{V \otimes W} = \chi_V \cdot \chi_W$.
\end{enumerate}
\end{proposition}

\begin{definition}[Inner product on class functions]
For $k = \mathbb{C}$, the \emph{inner product} on class functions is
\[
    \langle \alpha, \beta \rangle_G \coloneqq \frac{1}{|G|} \sum_{g \in G} \alpha(g)\,\overline{\beta(g)}.
\]
\end{definition}

\begin{theorem}[First Orthogonality Relation]
\label{thm:orth1}
Let $V_1,\ldots,V_r$ be the pairwise non-isomorphic irreducible representations
of a finite group $G$ over $\mathbb{C}$.  Then
\[
    \langle \chi_{V_i}, \chi_{V_j} \rangle_G = \delta_{ij}.
\]
In other words, the irreducible characters form an \emph{orthonormal} family in
$\mathrm{Cl}(G)$.
\end{theorem}

\begin{proof}
By Schur's Lemma (Corollary~\ref{cor:schur_decomp}) applied to
$\mathrm{Hom}_G(V_i, V_j)$, one computes
\[
    \langle \chi_{V_i}, \chi_{V_j} \rangle_G
    = \frac{1}{|G|} \sum_{g \in G} \chi_{V_i}(g)\,\overline{\chi_{V_j}(g)}
    = \dim_{\mathbb{C}} \mathrm{Hom}_G(V_i, V_j)
    = \delta_{ij}. \qedhere
\]
\end{proof}

\begin{theorem}[Second Orthogonality Relation]
\label{thm:orth2}
Let $C_1,\ldots,C_r$ be the conjugacy classes of $G$ and $g_s \in C_s$ any
representative.  Then
\[
    \sum_{i=1}^{r} \chi_{V_i}(g_s)\,\overline{\chi_{V_i}(g_t)} = \frac{|G|}{|C_s|}\,\delta_{st}.
\]
\end{theorem}

\begin{corollary}[Number of irreducibles]
The number of irreducible representations of a finite group equals the number of
its conjugacy classes.
\end{corollary}

\begin{corollary}[Decomposition formula]
\label{cor:decomp}
For any representation $(V,\rho)$ with character $\chi_V$, the multiplicity of
the irreducible $V_i$ in $V$ is
\[
    m_i = \langle \chi_V, \chi_{V_i} \rangle_G.
\]
\end{corollary}

\begin{example}[Character table of $S_3$]
The symmetric group $S_3$ has order $6$ and three conjugacy classes:
$C_1 = \{e\}$, $C_2 = \{(12),(13),(23)\}$, $C_3 = \{(123),(132)\}$.
The three irreducible representations are: the trivial representation $\mathbf{1}$,
the sign representation $\mathrm{sgn}$, and the standard $2$-dimensional
representation $\mathrm{std}$.

\begin{center}
\begin{tabular}{lccc}
\toprule
 & $e$ & $(12)$ & $(123)$ \\
\midrule
$\chi_{\mathbf{1}}$ & $1$ & $1$ & $1$ \\
$\chi_{\mathrm{sgn}}$ & $1$ & $-1$ & $1$ \\
$\chi_{\mathrm{std}}$ & $2$ & $0$ & $-1$ \\
\bottomrule
\end{tabular}
\end{center}

One verifies $\sum_i (\chi_{V_i}(e))^2 = 1 + 1 + 4 = 6 = |S_3|$, which is a
necessary identity for the complete list of irreducibles.
\end{example}

\begin{theorem}[Burnside's dimension formula]
\label{thm:burnside_dim}
Let $G$ be a finite group with irreducible complex representations $V_1,\ldots,V_r$
of dimensions $d_1,\ldots,d_r$.  Then
\[
    \sum_{i=1}^{r} d_i^2 = |G|.
\]
\end{theorem}

%% ============================================================
\section{Representations of $S_n$}
\label{sec:sn}
%% ============================================================

The symmetric group $S_n$ is one of the most important groups in mathematics and
physics, and its representation theory is completely governed by the combinatorics
of \emph{partitions} and \emph{Young tableaux}.

\begin{definition}
A \emph{partition} of $n$ is a weakly decreasing sequence of positive integers
$\lambda = (\lambda_1 \geq \lambda_2 \geq \cdots \geq \lambda_k > 0)$ with
$\lambda_1 + \cdots + \lambda_k = n$.  We write $\lambda \vdash n$.
\end{definition}

\begin{definition}
The \emph{Young diagram} of $\lambda \vdash n$ consists of $n$ boxes arranged in
rows from top to bottom, with $\lambda_i$ boxes in row $i$.  A \emph{Young
tableau} of shape $\lambda$ is a filling of the boxes with the numbers $1,\ldots,n$;
it is \emph{standard} if entries increase along each row from left to right and
down each column from top to bottom.
\end{definition}

\begin{example}
The partition $\lambda = (3,2) \vdash 5$ has Young diagram:
\[
\ytableaushort{{\phantom{1}}{\phantom{1}}{\phantom{1}},{\phantom{1}}{\phantom{1}}}
\quad\leftrightarrow\quad
\begin{array}{|c|c|c|}
\hline
\phantom{1} & \phantom{1} & \phantom{1} \\
\hline
\phantom{1} & \phantom{1} \\
\cline{1-2}
\end{array}
\]
The standard tableaux of this shape are those fillings with $\{1,2,3,4,5\}$
in which rows and columns are increasing; there are $5$ of them.
\end{example}

\begin{definition}[Specht module]
For $\lambda \vdash n$, the \emph{Specht module} $S^\lambda$ is a specific
irreducible $\mathbb{C}[S_n]$-module constructed as follows.  Fix a Young tableau
$T$ of shape $\lambda$.  Let
\[
    R(T) = \{\sigma \in S_n : \sigma \text{ preserves each row of } T\},
\quad
    C(T) = \{\sigma \in S_n : \sigma \text{ preserves each column of } T\}.
\]
Define the \emph{Young symmetrizer}
\[
    c_T = \underbrace{\left(\sum_{\sigma \in R(T)} \sigma\right)}_{a_T} \cdot
          \underbrace{\left(\sum_{\tau \in C(T)} \mathrm{sgn}(\tau)\,\tau\right)}_{b_T}.
\]
Then $S^\lambda = \mathbb{C}[S_n] \cdot c_T$ is an irreducible $\mathbb{C}[S_n]$-module,
independent of the choice of $T$ up to isomorphism.
\end{definition}

\begin{theorem}[Classification of irreducibles of $S_n$]
\label{thm:sn_class}
The modules $\{S^\lambda : \lambda \vdash n\}$ form a complete set of pairwise
non-isomorphic irreducible complex representations of $S_n$.  The dimension of
$S^\lambda$ equals the number of standard Young tableaux of shape $\lambda$,
given by the \emph{hook-length formula}:
\[
    \dim S^\lambda = \frac{n!}{\prod_{\square \in \lambda} h(\square)},
\]
where $h(\square)$ denotes the \emph{hook length} at box $\square$ (the number
of boxes directly to the right of or directly below $\square$, plus one for
$\square$ itself).
\end{theorem}

\begin{example}[Frobenius formula]
The character $\chi^\lambda$ of the Specht module $S^\lambda$ can be computed via
the Frobenius character formula.  For a permutation $\sigma \in S_n$ with cycle
type $\mu = (1^{m_1} 2^{m_2} \cdots)$,
\[
    \chi^\lambda(\sigma) = \langle s_\lambda, p_\mu \rangle,
\]
where $s_\lambda$ is the Schur polynomial, $p_\mu = \prod_i p_{\mu_i}$ is the
power-sum symmetric function, and the inner product is the Hall inner product on
symmetric functions.
\end{example}

\begin{example}[Irreducibles of $S_4$]
The partitions of $4$ are $(4)$, $(3,1)$, $(2,2)$, $(2,1,1)$, $(1,1,1,1)$,
giving $5$ irreducible representations of dimensions $1, 3, 2, 3, 1$ respectively.
One checks $1^2 + 3^2 + 2^2 + 3^2 + 1^2 = 1+9+4+9+1 = 24 = 4!$ as required.
\end{example}

%% ============================================================
\section{Representations of $\mathrm{GL}_n$ and $\mathrm{SL}_n$}
\label{sec:gln}
%% ============================================================

The theory of representations of compact Lie groups, or more generally of
reductive algebraic groups, is vastly richer than the finite-group case.  We
sketch the highest-weight classification for $\mathrm{GL}_n(\mathbb{C})$ and
$\mathrm{SL}_n(\mathbb{C})$.

\begin{definition}
A representation $(\rho, V)$ of $\mathrm{GL}_n(\mathbb{C})$ is
\emph{polynomial} (resp.\ \emph{rational}) if each matrix coefficient
$g \mapsto \langle v^*, \rho(g)v\rangle$ is a polynomial (resp.\ rational) function
of the entries of $g$.
\end{definition}

The irreducible polynomial representations of $\mathrm{GL}_n$ are indexed by
\emph{dominant weights}: sequences $\lambda = (\lambda_1 \geq \lambda_2 \geq
\cdots \geq \lambda_n \geq 0)$ of non-negative integers.  These correspond
precisely to partitions $\lambda$ with at most $n$ parts.

\begin{theorem}[Highest-weight classification]
\label{thm:hwt}
For each dominant weight $\lambda = (\lambda_1 \geq \cdots \geq \lambda_n \geq 0)$
there exists a unique (up to isomorphism) irreducible polynomial representation
$V_\lambda$ of $\mathrm{GL}_n(\mathbb{C})$.  Every irreducible polynomial
representation is of this form.  The representation $V_\lambda$ is precisely the
Schur module associated to the partition $\lambda$.
\end{theorem}

\begin{theorem}[Weyl's Character Formula]
\label{thm:weyl_char}
Let $G = \mathrm{SL}_n(\mathbb{C})$ with maximal torus $T$ of diagonal matrices
of determinant $1$.  For a dominant integral weight $\lambda$, the character of
the irreducible representation $V_\lambda$ evaluated on $\mathrm{diag}(t_1,\ldots,t_n)
\in T$ is
\[
    \chi_{V_\lambda}(t) = \frac{\sum_{w \in W} \epsilon(w)\, t^{w(\lambda+\rho)-\rho}}{\prod_{\alpha > 0}(1 - t^{-\alpha})},
\]
where $W = S_n$ is the Weyl group, $\epsilon(w) = \det(w)$, $\rho =
\tfrac{1}{2}\sum_{\alpha > 0}\alpha$ is the Weyl vector (half the sum of positive
roots), and the product runs over positive roots $\alpha$.
\end{theorem}

\begin{example}[Representations of $\mathrm{SL}_2$]
For $G = \mathrm{SL}_2(\mathbb{C})$, the dominant weights are non-negative
integers $n \geq 0$.  The irreducible representation $V_n$ of highest weight $n$
has dimension $n+1$ and can be realized as the space of homogeneous polynomials
of degree $n$ in two variables, with $\mathrm{SL}_2$ acting by substitution.
The character is
\[
    \chi_{V_n}\!\left(\begin{smallmatrix}t & 0 \\ 0 & t^{-1}\end{smallmatrix}\right)
    = t^n + t^{n-2} + \cdots + t^{-n} = \frac{t^{n+1} - t^{-(n+1)}}{t - t^{-1}}.
\]
These are the Chebyshev polynomials in disguise, and they coincide with the
spin-$n/2$ representations of $\mathrm{SU}(2)$ in quantum mechanics.
\end{example}

\begin{proposition}[Schur functors]
For a partition $\lambda \vdash d$ with at most $n$ parts, the Schur functor
$\mathbb{S}^\lambda$ associates to any vector space $E$ of dimension $n$ the
representation
\[
    \mathbb{S}^\lambda E \coloneqq (E^{\otimes d} \otimes S^\lambda)^{S_d},
\]
where $S^\lambda$ is the Specht module of $S_d$.  This gives the irreducible
$\mathrm{GL}(E)$-module $V_\lambda$, establishing a canonical link between the
representation theories of $S_d$ and $\mathrm{GL}_n$.
\end{proposition}

%% ============================================================
\section{Applications}
\label{sec:applications}
%% ============================================================

\subsection{Burnside's $pq$-Theorem}

\begin{theorem}[Burnside, 1904]
\label{thm:burnside_pq}
Let $G$ be a finite group of order $p^a q^b$ where $p$ and $q$ are prime and
$a, b \geq 0$.  Then $G$ is solvable.
\end{theorem}

\begin{proof}[Proof sketch]
The proof uses two character-theoretic lemmas.

\begin{lemma}
\label{lem:b1}
If $\chi$ is an irreducible character of $G$ of degree $d$ and $C$ is a
conjugacy class with $|C|$ coprime to $d$, then for any $g \in C$ either
$\chi(g) = 0$ or $\rho(g)$ is a scalar matrix.
\end{lemma}

\begin{lemma}
\label{lem:b2}
A finite group is solvable if and only if every non-trivial irreducible character
has degree divisible by some prime.
\end{lemma}

Using Sylow theory, one finds in $G$ a conjugacy class $C$ of size $p^s$ for some
$s \geq 1$.  Lemma~\ref{lem:b1} forces each irreducible character of degree not
divisible by $p$ to vanish on $C$ or to make $\rho(g)$ a scalar.  The column
orthogonality relation then forces $G$ to have a normal subgroup, contradicting
simplicity unless $G$ is abelian.  Inducting on $|G|$ yields solvability.
\end{proof}

\begin{remark}
Burnside's original 1904 proof was the first use of character theory to prove a
purely group-theoretic result.  A character-free proof was found by Goldschmidt
and Bender only in the 1970s, and it is considerably more involved.  This
illustrates the power of representation-theoretic methods.
\end{remark}

\subsection{Molecular Symmetry}

In chemistry, the symmetry group of a molecule (such as the rotation group of a
regular polyhedron) acts on the space of electron orbitals or vibrational modes.
The decomposition of these (typically reducible) representations into irreducibles
explains the degeneracy structure of energy levels and the selection rules for
spectroscopic transitions.

For example, the methane molecule $\mathrm{CH}_4$ has symmetry group $T_d \cong
S_4$.  The $2p$ orbitals of the four hydrogen atoms form the $4$-dimensional
permutation representation, which decomposes as $A_1 \oplus T_2$ (using
standard spectroscopic notation), accounting for the observed bonding pattern.

\subsection{Crystallography}

The $230$ crystallographic space groups in three dimensions classify all possible
periodic arrangements of atoms.  Each such group $\Gamma$ fits into a short exact
sequence $1 \to \mathbb{Z}^3 \to \Gamma \to P \to 1$, where $P$ is a
\emph{point group} (a finite subgroup of $\mathrm{O}(3)$).  The irreducible
representations of $P$ classify the symmetry types of phonon modes at the
$\Gamma$-point of the Brillouin zone, with direct applications to the
interpretation of Raman and infrared spectra.

%% ============================================================
\section{Connection to the Langlands Program}
\label{sec:langlands}
%% ============================================================

The Langlands program, initiated in Robert Langlands' 1967 letter to André Weil,
is a vast web of conjectures and theorems predicting deep correspondences between
three seemingly disparate objects:
\begin{enumerate}
  \item \textbf{Automorphic representations}: certain representations of adelic
        groups $G(\mathbb{A}_{\mathbb{Q}})$ that generalize modular forms and
        Dirichlet characters.
  \item \textbf{Galois representations}: continuous homomorphisms $\rho \colon
        \mathrm{Gal}(\bar{\mathbb{Q}}/\mathbb{Q}) \to {}^L G$ into the
        Langlands dual group.
  \item \textbf{L-functions}: meromorphic functions on $\mathbb{C}$ encoding
        arithmetic information via Euler products and satisfying functional equations.
\end{enumerate}

\subsection{Galois Representations}

Let $\ell$ be a prime.  An \emph{$\ell$-adic Galois representation} is a
continuous homomorphism
\[
    \rho \colon \mathrm{Gal}(\bar{\mathbb{Q}}/\mathbb{Q}) \longrightarrow \mathrm{GL}_n(\mathbb{Z}_\ell),
\]
where $\mathbb{Z}_\ell$ carries the $\ell$-adic topology and the Galois group
carries the Krull topology.  The representation theory developed in previous
sections applies here with $G = \mathrm{Gal}(\bar{\mathbb{Q}}/\mathbb{Q})$, a
profinite group.

\begin{example}[Cyclotomic character]
The \emph{cyclotomic character} $\chi_\ell \colon \mathrm{Gal}(\bar{\mathbb{Q}}/\mathbb{Q})
\to \mathbb{Z}_\ell^\times$ is defined by $\sigma(\zeta_{\ell^n}) = \zeta_{\ell^n}^{\chi_\ell(\sigma)}$
for all $\ell^n$-th roots of unity $\zeta_{\ell^n}$.  This is an irreducible
one-dimensional Galois representation encoding the arithmetic of cyclotomic fields.
\end{example}

\begin{example}[Galois representation from an elliptic curve]
Let $E/\mathbb{Q}$ be an elliptic curve.  The Tate module
\[
    T_\ell(E) = \varprojlim_n E[\ell^n](\bar{\mathbb{Q}}) \cong \mathbb{Z}_\ell^2
\]
carries a natural action of $\mathrm{Gal}(\bar{\mathbb{Q}}/\mathbb{Q})$, giving
a two-dimensional $\ell$-adic representation
$\rho_{E,\ell} \colon \mathrm{Gal}(\bar{\mathbb{Q}}/\mathbb{Q}) \to \mathrm{GL}_2(\mathbb{Z}_\ell)$.
The L-function $L(E,s) = L(\rho_{E,\ell}, s)$ is the same L-function that appears
in the Birch and Swinnerton-Dyer Conjecture.
\end{example}

\subsection{The Langlands Correspondence}

\begin{conjecture}[Global Langlands Correspondence for $\mathrm{GL}_n$]
\label{conj:langlands}
There is a natural bijection between:
\begin{enumerate}[label=\textup{(\roman*)}]
  \item Cuspidal automorphic representations $\pi$ of $\mathrm{GL}_n(\mathbb{A}_{\mathbb{Q}})$.
  \item Irreducible $n$-dimensional complex representations of the
        Weil--Deligne group $W_{\mathbb{Q}}' = W_{\mathbb{Q}} \ltimes \mathbb{C}$
        (or, conjecturally, of the Langlands group $\mathcal{L}_{\mathbb{Q}}$).
\end{enumerate}
Under this bijection, the L-function $L(s,\pi)$ on the automorphic side equals
the Artin L-function $L(s,\rho)$ on the Galois side.
\end{conjecture}

\begin{remark}
Conjecture~\ref{conj:langlands} remains open in general.  Major special cases
have been established:
\begin{itemize}
  \item $n=1$: Class field theory (Kronecker--Weber, Artin reciprocity).
  \item $G = \mathrm{GL}_2$, weight $2$ newforms: Taylor--Wiles (1995), as part of
        the proof of Fermat's Last Theorem.
  \item Function-field analogue: Drinfeld (1974, $n=2$) and Lafforgue (2002, general $n$).
  \item Geometrical Langlands over complex curves: Frenkel--Ben-Zvi and others.
\end{itemize}
The recent work of V.\ Lafforgue on the ``automorphic to Galois'' direction over
function fields, and of Scholze--Fargues using perfectoid geometry and the
Fargues--Fontaine curve, represent the current frontier.
\end{remark}

\subsection{Modularity and Fermat's Last Theorem}

The connection between representation theory and number theory becomes especially
visible in the proof of Fermat's Last Theorem.  Suppose $a^p + b^p = c^p$ with
$p \geq 5$ prime.  Frey (1986) associated to this hypothetical solution the
elliptic curve $E_{a,b,c}$, and Ribet (1990) proved that the associated Galois
representation $\rho_{E,p}$ would give rise to a weight-$2$ modular form of level
$2$.  But there are no such forms (by exhaustion of small levels), a contradiction.
The key step---that every semistable elliptic curve over $\mathbb{Q}$ is modular
(i.e.\ its Galois representation comes from a modular form)---was proved by
Wiles and Taylor--Wiles in 1995.  This modularity theorem is a landmark instance
of the Langlands correspondence in dimension $n = 2$.

%% ============================================================
\section*{References}
%% ============================================================

\begin{thebibliography}{99}

\bibitem{Serre77}
J.-P.\ Serre,
\textit{Linear Representations of Finite Groups},
Graduate Texts in Mathematics, vol.~42,
Springer-Verlag, New York, 1977.

\bibitem{FultonHarris91}
W.\ Fulton and J.\ Harris,
\textit{Representation Theory: A First Course},
Graduate Texts in Mathematics, vol.~129,
Springer-Verlag, New York, 1991.

\bibitem{JamesLiebeck01}
G.\ James and M.\ Liebeck,
\textit{Representations and Characters of Groups},
2nd ed., Cambridge University Press, Cambridge, 2001.

\bibitem{BumpLie04}
D.\ Bump,
\textit{Lie Groups},
Graduate Texts in Mathematics, vol.~225,
Springer-Verlag, New York, 2004.

\bibitem{GelfandManin03}
S.\ I.\ Gelfand and Yu.\ I.\ Manin,
\textit{Methods of Homological Algebra},
2nd ed., Springer Monographs in Mathematics,
Springer-Verlag, Berlin, 2003.

\bibitem{Langlands70}
R.\ P.\ Langlands,
\textit{Problems in the Theory of Automorphic Forms},
in: Lectures in Modern Analysis and Applications III,
Lecture Notes in Mathematics, vol.~170,
Springer-Verlag, Berlin, 1970, pp.~18--61.

\bibitem{TaylorWiles95}
R.\ Taylor and A.\ Wiles,
Ring theoretic properties of certain Hecke algebras,
\textit{Ann.\ of Math.} \textbf{141} (1995), 553--572.

\bibitem{Lafforgue02}
L.\ Lafforgue,
Chtoucas de Drinfeld et correspondance de Langlands,
\textit{Invent.\ Math.} \textbf{147} (2002), 1--241.

\bibitem{BumpSchilling17}
D.\ Bump and A.\ Schilling,
\textit{Crystal Bases: Representations and Combinatorics},
World Scientific, Singapore, 2017.

\bibitem{Neukirch99}
J.\ Neukirch, A.\ Schmidt and K.\ Wingberg,
\textit{Cohomology of Number Fields},
Grundlehren der mathematischen Wissenschaften, vol.~323,
Springer-Verlag, Berlin, 1999.

\end{thebibliography}

\end{document}
